\documentclass[10pt]{article}

\usepackage[utf8]{inputenc}

\usepackage[T1]{fontenc}

\usepackage{amsmath}

\usepackage{amsfonts}

\usepackage{amssymb}

\usepackage[version=4]{mhchem}

\usepackage{stmaryrd}

\usepackage{mathrsfs}

\usepackage{bbold}

\usepackage{graphicx}

\usepackage[export]{adjustbox}

\graphicspath{ {./images/} }



\begin{document}

ференциальных операторов бесконечного порядка), в теории гиперфункций, Фурье гиперфункций, Фурье ультрагиперфункций и т. Д.



Аналитический функционал в узком смысл е - элемент из $\mathscr{H}^{\prime}(O)$, то есть линейный непрерывный функционал на пространстве $\mathscr{H}(O)$ всех аналитич. функций в области $O \subset \mathbb{C}^{n}$ с топологией равномерной сходимости на компактах в $O$, задаваемой полунормами



\[

|\Phi|_{K}=\sup \{|\Phi(\zeta)| ; \zeta \in K\}, \Phi \in \mathscr{H}(O),

\]



где $K$ пробегает семейство всех компактов в $O$, так что $\mathscr{H}(O)$ - пространство Фреше - Шварца. Говорят, что А. ф. $g \in \mathscr{H}^{\prime}(O)$ сосредоточен на компакте $K \subset O$, если для любой ограниченной области $\Omega(K \subset \Omega$, замыкание $\bar{\Omega} \subset O)$ существует постоянная $C$ такая, что



\[

|(g, \Phi)| \leqslant C|\Phi|_{\bar{\Omega}}, \Phi \in \mathscr{H}(O) .

\]



Каждый А. ф. $g \in \mathscr{P}^{\prime}(O)$ сосредоточен на нек-ром компакте, лежащем в $O$, однако нельзя определить носитель $g$, поскольку не существует наименьшего компакта, на котором $g$ сосредоточен. Для всякого А. ф. $g \in \mathscr{H}^{\prime}(O)$ существует мера $\mu$ с компактным носителем supp $\mu \subset O$ такая, что



\[

(g, \Phi)=\int \Phi(z) d \mu(z), \Phi \in \mathscr{H}(O) .

\]



На $\mathscr{f}^{\prime}(O)$ определены Лапласа преобразование и свертка.



Вещественный аналитический функциона л в области $\omega \subset \mathbb{R}^{n}$ есть, по определению, элемент из $\mathcal{A}^{\prime}(\omega)$, то есть линейный непрерывный функционал на основном пространстве



\[

\mathscr{A}(\omega)=\operatorname{limind}_{\Omega \supset \omega} \mathscr{H}(\Omega), \quad,

\]



где индуктивный предел определенных выше пространств $\mathscr{H}(\Omega)$ берется по всем областям $\Omega \subset \mathbb{C}^{n}$, содержащим $\omega$. Для всякой области $\omega \subset \mathbb{R}^{n}$ имеется плотное непрерывное вложение $\mathscr{d}\left(\mathbb{R}^{n}\right) \subset \mathscr{d}(\omega)$, так что $\mathscr{d}^{\prime}\left(\mathbb{R}^{n}\right)-$ естественное объемлющее пространство всех вещественных А. Ф. Говорят, что вещественный А. $\phi . g \in d^{\prime}(\omega)$ сосредоточен на вещественном компакте $K \subset \omega$, если для всякой ограниченной области $(\Omega) \subset \mathbb{C}^{n}$, содержащей $K$, существует постоянная $C$ такая, что



\[

|(g, \Phi)| \leqslant C|\Phi|_{\bar{\Omega}}, \Phi \in \mathscr{H}\left(\mathbb{C}^{n}\right) .

\]



Каждый вещественный А. Ф. сосредоточен на нек-ром вещественном компакте, определен носитель supp $g$ всякого вещественного А. Ф. $g$ как наименьший вещественный компакт, на к-ром он сосредоточен. Множество вещественных А. ф., носители к-рых лежат в компакте $K \subset \mathbb{R}^{n}$, отождествляется с сопряженным пространством $\mathscr{d}^{\prime}[K]$, где



\[

\mathscr{A}[K]=\underset{\Omega \supset K}{\operatorname{limind}} \mathscr{H}(\Omega),

\]



индуктивный предел берется по всем комплексным открытым множествам $\Omega$, содержащим $K$. Для вещественных А. ф. определены преобразование Лапласа и свертка.



Лит.: [1] Хе рмандер Л., Введение в теорию функций нескольких комплексных переменных, пер. с англ., М., 1968; [2] Шап и ра П., Теория гиперфункций, пер. с франц., М., 1972; [3] Еф им о в Г. В., Нелокальные взаимодействия квантованных полей, М., $\begin{array}{ll}1977 . & \text { B. В. Жаринов. }\end{array}$ АНАЛИТИЧЕСКИЙ ЭЛЕМЕНТ, эл е ме нт ан ал и ти ч е с ко й фу н к ц и и,- см. Полная аналитическая функция. АНАЛИТИЧЕСКОЕ ПРЕДСТАВЛЕНИЕ О б О бЩ Н н Н ОЙ Функци и - представление обобщенной функции в виде суммы обобщенных граничных значений аналитических функший. Широко применяется в теории обобщенных функций, в теории линейных уравнений с частными производными, в аксиоматич. квантовой теории поля, явилось мотивировкой введения понятия гиперфункции как дальнейшего расширения классич. понятия функции.



Пусть $O$ - область в $\mathbb{R}^{n}$ и $\mathscr{\mathscr { F }}(O)$ - основное пространство, состоящее из функций, определенных в области $O$. Аналитическим представлением обоб- ще н ной функци и $f \in \mathscr{F}^{\prime}(O)$ называется всякое равенство вида



\[

f=\sum_{j=1}^{N} \mathrm{bv} f_{j}

\]



где справа стоят обобщенные граничные значения bv $f_{j}$ в $\mathscr{F}^{\prime}(O)$ функций $f_{j}(z)$, аналитических в тубоидах $D_{j}(O)$. Во всех разумных пространствах всякая обобщенная функция обладает А. П. Неоднозначность А. П. обобщенной функции описывается Боголюбова теоремой «острие клина», точнее, ее вариантом:



\[

\sum_{j=1}^{N} \mathrm{bv} f_{j}=0

\]



тогда и только тогда, когда существуют функции $f_{j k}(z)=$ $=-f_{k j}(z)$, аналитические в тубоидах $D_{j k}(O)=D_{k j}(O) \supset$ $\supset D_{j}(O) \cup D_{k}(O), j, k=1, \ldots, N$, имеющие обобщенные граничные значения в $\widetilde{\mathscr{T}}^{\prime}(O)$, и такие, что



\[

f_{j}(z)=\sum_{k=1}^{N} f_{j k}(z), \quad z \in D_{j}(O), \quad j=1, \ldots, N .

\]



Приведенные результаты справедливы, напр., в классах Швариа распределений, ультрараспределений, Фурье гипер$\begin{array}{ll}\text { функций и Фурье ультрагиперфункций. } & \text { В. В. Жаринов. }\end{array}$ АНАЛИТИЧЕСКОЕ ПРОДОЛЖЕНИЕ ФУ НКЦИ и - Доопределение функции $f_{0}(z)$, голоморфной в области $D_{0}$ плоскости $\mathbb{C}$ комплексного переменного $z$, то есть однозначной и аналитической в $D_{0}$, до функции $f(z)$, голоморфной в более обширной области $D$, содержащей $D_{0}$, причем в исходной области $D_{0}$ функция $f(z)$ совпадает с $f_{0}(z)$. Отправным понятием в теории А. п. является понятие а на литического элемента, или просто элемента $\left(D_{0}, f_{0}\right)$, состоящего из области $D_{0}$ и голоморфной в $D_{0}$ функции $f_{0}$. При этом теоретически всегда возможно ограничиться рассмотрением простейших канонических эл е м е н т о в $\left(K_{0}, f_{0}\right)$, состоящих из степенного ряда



\[

f_{0}=f_{0}(z)=\sum_{n=0}^{\infty} a_{n}\left(z-z_{0}\right)^{n}

\]



и его круга сходимости $K_{0}=\left\{z:\left|z-z_{0}\right|<R_{0}\right\}$ с центром $z_{0} \in \mathbb{C}$ и радиусом $R_{0}>0$. Приняв за центр ряда другую точку $z_{1} \neq z_{0}$ этого круга и пересчитав коэффициенты, напр. по формулам Тейлора $b_{n}=f^{(n)}\left(z_{1}\right) / n$ !, получают новый элемент $\left(K_{1}, f_{1}\right)$, состоящий из степенного ряда



\[

f_{1}=f_{1}(z)=\sum_{n=0}^{\infty} b_{n}\left(z-z_{1}\right)^{n}

\]



и его круга сходимости $K_{1}=\left\{z:\left|z-z_{1}\right|<R_{1}\right\}$. Может оказаться (см. рис. 1), что $R_{1}>R_{0}-\left|z_{1}-z_{0}\right|$, то есть круг $K_{1}$



Рис. 1. Непосредственное аналитическое продолжение.



\begin{center}

\includegraphics[max width=\textwidth]{2023_07_08_7a6489a9de87fcc48e3dg-1}

\end{center}



выходит за пределы круга $K_{0}$. Тогда говорят, что канонич. элементы $\left(K_{1}, f_{1}\right)$ и $\left(K_{0}, f_{0}\right)$ составляют н е п о с р е д с т в е н н о е А. п. друг друга через общую (заштрихованную) часть кругов $K_{0}$ и $K_{1}$. Элемент $\left(K_{0}, f_{0}\right)$ допускает А. П. в граничную точку $\zeta$ его круга сходимости, $\left|\zeta-z_{0}\right|=R_{0}$, если для нек-рого непосредственного А. П. $\left(K_{1}, f_{1}\right)$ эта точка $\zeta$ попадает в



АНАЛИТИЧЕСКОЕ 25





\end{document}